\documentclass[11pt]{article}
\usepackage[margin=1in]{geometry}
\usepackage{amsmath,amssymb}
\usepackage{booktabs}

\title{Disproof of Conjecture 00000001092:\\Small Complete Caps in $PG(2,27)$}
\author{}
\date{}

\begin{document}
\maketitle

\section{The conjecture}

The original statement reads, with its definition sentence quoted verbatim:

\begin{quote}
``Definition: A complete cap of PG(2,27) is a maximal set with no three collinear points. Conjecture: Its minimum size is 21, with exactly two isomorphism classes attaining this value; the second-smallest complete cap has size 22.''
\end{quote}

\section{The counterexample}

Work over $\mathbb{F}_{27}=\mathbb{F}_3[\alpha]$ with $\alpha^3=\alpha+1$; the polynomial $t^3-t+2$ has no root in $\mathbb{F}_3$ (evaluating at $0,1,2$ gives $2,2,2$), hence is irreducible, so this presentation is valid. An element is written $c_0+c_1\alpha+c_2\alpha^2$ and encoded by the digit triple $(c_0,c_1,c_2)$.

The plane $PG(2,27)$ has $27^2+27+1=757$ points, represented by the canonical triples $(x:y:z)$ whose first nonzero coordinate is $1$; three points are collinear exactly when the determinant of their coordinate matrix vanishes in $\mathbb{F}_{27}$.

Table~\ref{tab:cap} lists $14$ canonical points. Direct computation over $\mathbb{F}_{27}$ (deterministic, fully reproducible) establishes:

\begin{enumerate}
\item \textbf{Cap:} no three of the $14$ points are collinear (all $\binom{14}{3}=364$ determinant checks are nonzero).
\item \textbf{Complete:} each of the remaining $743$ points of $PG(2,27)$ lies on a secant of the set, i.e., on a line through two of its points. Hence no point can be added without creating three collinear points, so the set is a maximal set with no three collinear points --- a complete cap in the sense of the definition above.
\end{enumerate}

\begin{table}[h]
\centering
\begin{tabular}{clll}
\toprule
$i$ & $x_i$ & $y_i$ & $z_i$\\
\midrule
$P_1$ & $(1)$ & $(2+\alpha+\alpha^2)$ & $(2+2\alpha+2\alpha^2)$\\
$P_2$ & $(1)$ & $(1)$ & $(2+\alpha+\alpha^2)$\\
$P_3$ & $(1)$ & $(2+\alpha^2)$ & $(0)$\\
$P_4$ & $(0)$ & $(1)$ & $(1+2\alpha+\alpha^2)$\\
$P_5$ & $(1)$ & $(1+\alpha)$ & $(\alpha+2\alpha^2)$\\
$P_6$ & $(1)$ & $(2+2\alpha^2)$ & $(1+2\alpha)$\\
$P_7$ & $(1)$ & $(2+\alpha+\alpha^2)$ & $(1+\alpha)$\\
$P_8$ & $(1)$ & $(1+\alpha^2)$ & $(2\alpha)$\\
$P_9$ & $(1)$ & $(1+\alpha+\alpha^2)$ & $(1+2\alpha+\alpha^2)$\\
$P_{10}$ & $(1)$ & $(2+\alpha^2)$ & $(1+2\alpha+\alpha^2)$\\
$P_{11}$ & $(1)$ & $(\alpha)$ & $(2+\alpha)$\\
$P_{12}$ & $(1)$ & $(1+2\alpha+\alpha^2)$ & $(1+\alpha)$\\
$P_{13}$ & $(1)$ & $(1)$ & $(2+2\alpha^2)$\\
$P_{14}$ & $(1)$ & $(0)$ & $(\alpha+\alpha^2)$\\
\bottomrule
\end{tabular}
\caption{A complete cap of $PG(2,27)$ of size $14$; each coordinate is an element of $\mathbb{F}_{27}=\mathbb{F}_3[\alpha]$, $\alpha^3=\alpha+1$, written as $c_0+c_1\alpha+c_2\alpha^2$.}
\label{tab:cap}
\end{table}

\section{Consequence for the conjecture}

A complete cap of size $14$ exists in $PG(2,27)$, so
\[
m(27) \mathrel{\mathop:}= \min\{|C| : C \text{ complete cap of } PG(2,27)\} \le 14 < 21.
\]
The leading claim ``its minimum size is $21$'' is therefore false. Furthermore, an independent randomized greedy search (40 restarts, each run exhausted over all $757$ points and hence producing a \emph{complete} cap by construction) yielded complete caps of sizes $14,15,16,17$ only; in particular complete caps of sizes $14$ and $15$ both exist, so the second-smallest complete cap has size at most $15 < 22$, and the residual claims (exactly two isomorphism classes at the minimum; second-smallest of size $22$) collapse as well. Conjecture 00000001092 is disproved in full.

\section{Machine verification}

Two independent certificates are included in this submission:

\begin{itemize}
\item \texttt{reproduce.py}: a deterministic, stdlib-only Python script that rebuilds $\mathbb{F}_{27}$ and the $757$ canonical points of $PG(2,27)$, then verifies that the $14$ listed points are canonical, distinct, pairwise non-collinear in triples, and that every outside point lies on a secant. It prints \texttt{PASS}.
\item \texttt{lean4/Main.lean}: a Lean~4 (v4.33.1, no Mathlib, no \texttt{native\_decide}) certificate in which
\begin{verbatim}
theorem plane_size       : plane.length = 757                    := by decide
theorem cap_small        : cap.length = 14 /\ cap.length < 21    := by decide
theorem cap_is_cap       : noThreeCollinear = true               := by decide
theorem cap_is_complete  : isComplete = true                     := by decide
theorem main_disproof    : disproof_claim = true                 := by decide
\end{verbatim}
are proved purely by kernel reduction. \texttt{Check.lean} runs \texttt{\#print axioms} on all five theorems; each reports ``does not depend on any axioms'' (zero axioms, zero \texttt{sorry}).
\end{itemize}

\end{document}
